\documentclass[11pt]{article}
\usepackage[T1]{fontenc}
\usepackage{lmodern,amsmath,amssymb,amsthm}
\usepackage[letterpaper,margin=0.95in]{geometry}
\usepackage[hidelinks]{hyperref}
\hypersetup{pdftitle={Optimal periodic resetting for a Gaussian target on the line},
 pdfauthor={Mathematical Physics Notes}}
\newtheorem{theorem}{Theorem}
\newtheorem{lemma}[theorem]{Lemma}
\newtheorem{proposition}[theorem]{Proposition}
\newcommand{\E}{\mathbb E}
\newcommand{\Prob}{\mathbb P}
\newcommand{\erfc}{\operatorname{erfc}}
\newcommand{\dd}{\,\mathrm d}
\title{Optimal periodic resetting for a Gaussian target on the line}
\author{Mathematical Physics Notes}
\date{Working paper \(\cdot\) 26 November 2026}
\begin{document}
\maketitle
\begin{abstract}
A Brownian searcher returns instantaneously to the origin at successive
deadlines until it hits a target. The target is sampled once from a centered
Gaussian distribution and remains fixed after each reset. We prove that a
constant deadline minimizes the mean search time over all positive, almost
surely finite, independent and identically distributed deadlines, including
laws with infinite mean. On a compact interval we select a deterministic
period by matching the spatial logarithmic slopes of the single-attempt
success probabilities. Two one-crossing properties give a global comparison
by convexity. Exhausting the half-line yields the full comparison and
attainment. The resulting period and mean time coincide with the previously
computed optimum within the periodic family.
\end{abstract}

\section{Model and result}
Let $B$ be standard one-dimensional Brownian motion, $\E B_t^2=t$.
Sample $Z\sim N(0,1)$ independently at the start of the search. Each attempt
starts at zero with a fresh Brownian path. At its deadline the searcher
returns instantaneously to zero if it has not hit $Z$. The deadlines have
a common law $\mu$ on $(0,\infty)$ and are mutually independent and
independent of $Z$ and of all paths. Write $J(\mu)$ for the mean time of the
first hit, averaged over all these random variables, and $\delta_t$ for the
constant deadline $t$. Conditional on $X=|Z|=x$, reflection symmetry allows
us to put the target at $+x$. Thus $p(x)=\sqrt{2/\pi}\,e^{-x^2/2}$ for $x>0$.

\begin{theorem}\label{th:main}
For every $\mu$ with $J(\mu)<\infty$ there is a $t>0$ such that
$J(\delta_t)\le J(\mu)$. Moreover,
\[
       \inf_{\mu}J(\mu)=\min_{t>0}J(\delta_t)<\infty .
\]
\end{theorem}

Pinsky computed the optimum over periodic instantaneous resets and compared
it with Poissonian resets and Brownian bridge resets~\cite[Theorem~3]{pinsky}.
Theorem~\ref{th:main} upgrades that periodic optimum to a comparison with
every renewal law in the model. In particular a single period suffices when
the realized target distance is hidden. Policies using observations,
the attempt number, or the absolute time are outside the theorem.

\section{The renewal objective}
Write $H_x$ for the first time a fresh path hits $x$, and let $S\sim\mu$.
For a deterministic deadline the reflection principle gives
\[
 q_t(x)=\Prob(H_x<t)=\erfc\!\bigl(x/\sqrt{2t}\bigr),\qquad
 g_t(x)=-q_t'(x)=\sqrt{2/(\pi t)}\,e^{-x^2/(2t)} .
\]
Set $q=\E q_S$, $g=\E g_S$ and $a=\E\min(H_x,S)$, with the dependence of
$a$ on $x$ understood. Then $q>0$ and $g=-q'$ on $(0,\infty)$.
Differentiation under the expectation on compact subintervals away from zero
follows because $s^{-k} e^{-x^2/(2s)}$ is uniformly bounded there for
every $k>0$. The geometric sum over attempts at fixed $x$ gives mean
$a(x)/q(x)$, also when $a(x)=\infty$.

We record the kernel formula used in the companion note~\cite{note}.
For $a_t(x)=\int_0^t(1-q_v(x))\,dv$, the heat equation gives
$-a_t''=2q_t$, with $a_t(0)=0$ and $a_t'(\infty)=0$. The last condition
follows from $a_t'=\int_0^t g_v\,dv$ by dominated convergence; in the
interior the Gaussian factor permits the two differentiations.
Integrating twice and averaging over $S$ by Tonelli's theorem therefore gives
\begin{equation}\label{eq:J}
 a(x)=2\int_0^\infty\min(x,y)q(y)\,dy,\qquad
 J(\mu)=2\iint_{(0,\infty)^2}\min(x,y)p(x)\frac{q(y)}{q(x)}\,dx\,dy .
\end{equation}
Fix $I=[a,b]\subset(0,\infty)$ with $a<b$. Write $J_I(q)$ for the last
integral restricted to $I^2$. This is a truncation of an integral; the
Brownian paths still move on the whole line.

\section{Matching the slope of a mixture}
Write $h=-q'/q$ and $h_t=-q_t'/q_t$. These are logarithmic slopes in
\emph{distance}. The next lemma depends only on the Brownian search process.
\begin{lemma}\label{lem:match}
For every $u>0$ there is a unique $t=t(u)$ with $h_t(u)=h(u)$, and
$t(u)$ is continuous. For this fixed $t$,
\[
 h(x)-h_t(x)\ge0\quad(0<x<u),\qquad
 h(x)-h_t(x)\le0\quad(x>u).
\]
\end{lemma}
\begin{proof}
The identity
\begin{equation}\label{eq:mills}
 h_t(x)^{-1}=\frac{q_t(x)}{g_t(x)}
 =\int_0^\infty e^{-(2xv+v^2)/(2t)}\,dv
\end{equation}
shows that $1/h_t(u)$ increases strictly and continuously from zero to
infinity as $t$ increases. Since $h(u)>0$ is continuous, so is the
unique solution: periods on either side of $t(u)$ still bracket the root
for all sufficiently close $u'$.

For the matched $t$, consider the \emph{density} ratio
\[
 d(x)=g(x)/g_t(x)
 =\E\!\left[\sqrt{t/S}\,\exp\{(1/t-1/S)x^2/2\}\right].
\]
As a function of $x^2$ it is strictly convex unless $S=t$ almost surely,
in which case the assertion is immediate. Its derivatives on compact subsets
of $(0,\infty)$ are justified by the same Gaussian bound as for $g$. Consequently $d$ is decreasing
then increasing in $x$, with the possibility that one of the two parts is
empty. Define
\[
 \Delta(x)=q(x)-q_t(x)d(x)
          =\int_x^\infty g_t(y)[d(y)-d(x)]\,dy .
\]
On a decreasing part $\Delta'=-q_t d'>0$. Wherever $d$ is increasing
for all subsequent $x$, the integral is strictly positive. If $d$ is
strictly decreasing throughout, the integral is strictly negative.
Thus any zero of $\Delta$ is a crossing from negative to positive.
Now $(q/q_t)'=g_t\Delta/q_t^2$ and $\Delta(u)=0$ by matching. Taking the
negative logarithmic derivative of $q/q_t$ proves the two signs.
\end{proof}

\section{A global tangent inequality on a compact interval}
For $x<y$ let $A=\int_x^y h(v)\,dv$. Pair the contributions of $(x,y)$
and $(y,x)$ in~\eqref{eq:J}. The result is $2x$ times
$p(x)e^{-A}+p(y)e^A$, a convex function of $A$.
Take its tangent at $A_t=\int_x^y h_t(v)\,dv$, and integrate over all
$a<x<y<b$. By changing the order of integration we obtain
\begin{equation}\label{eq:tangent}
 J_I(q)-J_I(q_t)\ \ge\
 \int_a^b W_t(u)[h(u)-h_t(u)]\,du .
\end{equation}
All functions in this step are bounded on $I$. To determine the sign of
$W_t$, put $r=p/q_t^2$ and define the masses and averages
\begin{align*}
 U(u)&=\int_a^u xq_t(x)\,dx,&
 L(u)&=U(u)^{-1}\int_a^u xq_t(x)r(x)\,dx,\\
 V(u)&=\int_u^b q_t(y)\,dy,&
 R(u)&=V(u)^{-1}\int_u^b q_t(y)r(y)\,dy .
\end{align*}
In the tangent coefficient the derivative of the paired integrand is
$2xq_t(x)q_t(y)[r(y)-r(x)]$. Integrating over $x<u<y$ gives
\[
 W_t(u)=2U(u)V(u)G_t(u),\qquad G_t(u)=R(u)-L(u).
\]
Use the continuous limits $L(a)=r(a)$ and $R(b)=r(b)$.

\begin{lemma}\label{lem:weight}
For a fixed $t$, every interior zero of $G_t$ is a strict downward crossing.
If $G_t(a)=0$ it is negative immediately to the right; if $G_t(b)=0$ it is
positive immediately to the left. In particular $G_t$ has at most one zero
on $I$.
\end{lemma}
\begin{proof}
The Gaussian target gives $(\log r)'=2h_t-x$. Upon substituting $w=xv$
in~\eqref{eq:mills} we find
\[
 x/h_t(x)=\int_0^\infty e^{-w/t-w^2/(2tx^2)}\,dw .
\]
This increases strictly with $x$, so $(\log r)'/x$ decreases strictly.
On $I$ the function $r$ therefore either is strictly monotone or has one
strict interior maximum, increasing before it and decreasing after it.

At an interior zero $u$ write $m=L(u)=R(u)$. On the increasing side the
left average is strictly less than $r(u)$; on the decreasing side the
right average is strictly less. This also covers the monotone cases and the
maximum itself. Hence $m<r(u)$. The identities
\[
 L'(u)=\frac{u q_t(u)}{U(u)}[r(u)-L(u)],\qquad
 R'(u)=\frac{q_t(u)}{V(u)}[R(u)-r(u)]
\]
now give $G_t'(u)<0$.

If $G_t(a)=0$, then $r(a)$ equals the right average over $I$. We must
have $r'(a)>0$: otherwise the strict decrease of $(\log r)'/x$ makes $r$
strictly decreasing on $(a,b)$, so its average is below $r(a)$.
For $v\downarrow0$, the positive smooth left weight $xq_t(x)$ gives
$L(a+v)=r(a)+r'(a)v/2+o(v)$, by expanding numerator and denominator
(or averaging $r(x)-r(a)$ on the short interval). The derivative formula
for $R$ at $a$ gives $R'_+(a)=0$. Thus $G'_{t,+}(a)=-r'(a)/2<0$.
If $G_t(b)=0$, the same strict decrease of $(\log r)'/x$ forces $r'(b)<0$;
otherwise the left average would be below $r(b)$. Here $L'_-(b)=0$ and
$R(b-v)=r(b)-r'(b)v/2+o(v)$, giving $G'_{t,-}(b)=r'(b)/2<0$.
Continuity precludes two zeros since a second downward crossing would
require a return from negative to nonnegative between them.
\end{proof}

\begin{proposition}\label{prop:compact}
For each $I$ and admissible $\mu$ there exists a single $t>0$ such that
$J_I(q_t)\le J_I(q)$.
\end{proposition}
\begin{proof}
Let $F(u)=G_{t(u)}(u)$ with $t(u)$ as in Lemma~\ref{lem:match}. The
averages defining $G_t(u)$ are jointly continuous in $(t,u)$ up to the
endpoints for $t>0$: a shrinking-interval average tends to its endpoint
value uniformly for $t$ in a compact neighborhood. Thus $F$ is continuous.
If $F(a)\le0$, fix $t=t(a)$. Lemma~\ref{lem:weight} gives $G_t\le0$
throughout $I$ and Lemma~\ref{lem:match} gives $h-h_t\le0$. If $F(b)\ge0$,
$t=t(b)$ instead makes both functions nonnegative. In the remaining case
$F(a)>0>F(b)$. Choose an interior zero $u_0$ of $F$ and fix $t=t(u_0)$.
For this \emph{single} period $G_t$ and $h-h_t$ are both nonnegative on
the left of $u_0$ and nonpositive on the right. In every case the
right-hand side of~\eqref{eq:tangent} is nonnegative.
\end{proof}

\section{Passage to the half-line and attainment}
\begin{proof}[Proof of Theorem~\ref{th:main}]
For $n\ge2$ take $I_n=[1/n,n]$ and let $t_n$ be the comparison period
from Proposition~\ref{prop:compact}. Then
$J_{I_n}(q_{t_n})\le J_{I_n}(q)\le J(\mu)<\infty$.
Extend the nonnegative integrands by zero off $I_n^2$.
If a subsequence had $t_n\to\infty$, its ratio $q_{t_n}(y)/q_{t_n}(x)$
would converge to one. The resulting integral in~\eqref{eq:J} diverges
because $\int_x^\infty \min(x,y)\,dy=\infty$ for each $x>0$.
If $t_n\to0$, the Gaussian tail asymptotic
$q_t(x)\sim\sqrt{2t/\pi}\,x^{-1}e^{-x^2/(2t)}$ for fixed $x>0$
gives $q_{t_n}(y)/q_{t_n}(x)\to\infty$ on the positive-measure set $0<y<x$.
Fatou's lemma contradicts the finite bound in both cases.
Choose $n_j\to\infty$ such that $t_{n_j}\to t_0\in(0,\infty)$. Fatou again
gives $J(\delta_{t_0})\le\liminf_{j\to\infty} J_{I_{n_j}}(q_{t_{n_j}})\le J(\mu)$.

It remains to show that the periodic infimum is attained.
A finite competitor is $t=2$. Indeed, $a_2\le2$ and the same Gaussian tail
asymptotic, now at $x\to\infty$ with $t=2$, gives
$p(x)/q_2(x)=O(xe^{-x^2/4})$; the ratio is bounded near zero.
Take a minimizing sequence of periods with bounded costs.
Applying the preceding two Fatou arguments to the full integrals rules
out subsequences tending to zero or infinity. A subsequence converges to
a finite positive period and Fatou bounds its cost by the infimum. Since
constant clocks belong to the admissible class, the two infima coincide.
\end{proof}

\section{Scale, benchmark and relation to earlier results}
If the Brownian variance is $\nu t$ and the centered target variance is
$\sigma^2$, Brownian scaling multiplies both times in the unit model by
$\sigma^2/\nu$. In the convention with generator $D\partial_x^2$ one has
$\nu=2D$. For numerical evaluation within the periodic family,
integrating the Brownian survival probability gives
\[
 a_t(x)=t[1-q_t(x)]+x\sqrt{2t/\pi}\,e^{-x^2/(2t)}-x^2q_t(x).
\]
Using $J(\delta_t)=\int_0^\infty p(x)a_t(x)/q_t(x)\,dx$ gives a period
near $2.81592$ and mean time near $3.34972$ in the unit model, consistent
with Pinsky's values $2.82$ and $3.35$~\cite[Eq.~(1.13)]{pinsky}. The
accompanying Python script reproduces this numerical benchmark. Attainment
and global optimality follow from the comparison proof; uniqueness of the
minimizing period is a separate question.

Pal and Reuveni~\cite{pal} show why a deterministic deadline suffices for
a fixed completion-time law renewed independently at each restart. Here,
conditional on $X=x$, the relevant law is fixed, but the best period may
depend on $x$. Averaging those pointwise minima allows knowledge of $X$.
In the unconditional search, failures carry information about the original
target and the same $X$ enters the success probability of every attempt.
The comparison above selects a common period.

Evans and Ray~\cite{evans} studied precisely this distinction between fixed
and randomly located targets. They derived a conjugate target density
proportional to the square of the one-attempt success probability and
discussed the advantage of periodic resetting for narrow targets.
For a periodic clock the conjugate density is proportional to
$\erfc^2(x/\sqrt{2t})$ in our normalization. It has a Gaussian-type tail,
whereas the exact Gaussian distance density used here is proportional to
$e^{-x^2/2}$. The squared-success equality case from the companion
bound~\cite{note} also leaves this instance open: it would require
$q(x)=e^{-x^2/4}$, but every $q_t$ and every positive mixture of them is
convex on $(0,\infty)$, while $e^{-x^2/4}$ is concave near zero.
The compact comparison supplies the missing inequality for this target.

\paragraph{Manuscript status.}
This working paper has not undergone peer review. OpenAI Codex assisted
with the derivation, drafting and numerical check; a separate AI reading
examined the sign and endpoint arguments. The benchmark values are
reported for reproducibility. Priority for the global comparison remains
subject to independent literature review.

\begin{thebibliography}{4}
\bibitem{pinsky} R.~G. Pinsky, Comparison of Brownian jump and Brownian
bridge resetting in search for Gaussian target on the line and in space,
\emph{J. Phys. A: Math. Theor.} \textbf{56}, 335004 (2023).
\href{https://doi.org/10.1088/1751-8121/ace865}{doi:10.1088/1751-8121/ace865}.
\bibitem{pal} A. Pal and S. Reuveni, First Passage under Restart,
\emph{Phys. Rev. Lett.} \textbf{118}, 030603 (2017).
\href{https://doi.org/10.1103/PhysRevLett.118.030603}{doi:10.1103/PhysRevLett.118.030603}.
\bibitem{evans} M.~R. Evans and S. Ray, Stochastic Resetting Prevails Over
Sharp Restart for Broad Target Distributions, \emph{Phys. Rev. Lett.}
\textbf{134}, 247102 (2025).
\href{https://doi.org/10.1103/2wpq-gl2j}{doi:10.1103/2wpq-gl2j}.
\bibitem{note} Mathematical Physics Notes, A sharp bound for Brownian
search with renewal resetting, working paper (2026).
\href{https://github.com/daemonx666/math-physics-notes/blob/92bc9bc972a7755d4cc9ff2c768b29f5155a4621/papers/brownian-resetting/paper.pdf}{Repository copy}.
\end{thebibliography}
\end{document}
